% In this file you should put all LaTeX macros and settings to be used both by
% the pdf version and the web version.
% This should be most of your macros.

% The theorem-like environments defined below are those that appear by default
% in the dependency graph. See the README of leanblueprint if you need help to
% customize this. 
% The configuration below use the theorem counter for all those environments
% (this is what the [theorem] arguments mean) and never resets it.
% If you want for instance to number them within chapters then you can add
% [chapter] at the end of the next line.
\newtheorem{theorem}{Theorem}
\newtheorem{proposition}[theorem]{Proposition}
\newtheorem{lemma}[theorem]{Lemma}
\newtheorem{corollary}[theorem]{Corollary}

\theoremstyle{definition}
\newtheorem{definition}[theorem]{Definition}

% --- T51: the classical-input / assumption environment -------------------
% Registered per TASK_BOARD.md §1 (epistemic three-way split) and the
% blueprint/README.md T51 authoring note: interface fields of `ClassicalInputs`
% (citable classical theorems: Rubin, Mazur, Greenberg, MTT, Bannai–Kobayashi,
% de Shalit, Schneider/Perrin-Riou, Hasse) are rendered in this environment
% rather than as `lemma`/`proposition`/`theorem`, so a reader can see at a
% glance, by the header word alone ("Assumption" vs "Lemma"/"Theorem"), which
% statements are asserted on citation and which are Lean-kernel-proved.
% `\leanok` is never applied to a node in this environment (TASK_BOARD.md §1:
% these are the human-refereeable trust boundary, not proof obligations).
%
% NOTE for the graph-owning agent (not this task's fence): plasTeX's
% `plastexdepgraph` plugin only draws dependency-graph nodes for theorem
% *kinds* listed in the `thms=` option passed to
% `\usepackage[showmore, dep_graph]{blueprint}` in `src/web.tex` (default
% `definition+lemma+proposition+theorem+corollary`). `src/web.tex` is outside
% this task's fence (T51), so `assumption` nodes below render with their own
% distinct header/style in the PDF and HTML *prose* but are **not** picked up
% by the auto-generated dependency graph. The handful of assumption-flavoured
% nodes that a proved theorem's dependency chain must visibly pass through in
% the graph (`lem:comparison`, `eq:padicbsd`, `def:deltaE`, `prop:jetformula`,
% `hyp:sinnott`, `conj:EK`) are therefore instead declared as `definition`
% (one of the five default graph-visible kinds, box-shaped) and deliberately
% never carry `\leanok`, so they still read as visually distinct (box, no
% green fill) from proved `theorem`/`lemma`/`proposition`/`corollary` nodes
% (ellipse, green once proved). See blueprint/README.md for the one-line
% `src/web.tex` change (`thms=...+assumption`) that would let `assumption`
% itself join the graph, if a future task takes on that file.
\newtheorem{assumption}[theorem]{Assumption}

%% --- Paper macros, ported verbatim from main.tex (T51 / Wave 7 fix, 2026-08-18).
%% The blueprint restates the paper's results, so it needs the paper's notation.
%% Keep this block in sync with main.tex's preamble (main.tex lines 11-31).
%%
%% BRACE-TRANSPARENCY FIX (2026-08-18, W7-A): this blueprint builds with
%% xelatex + unicode-math (see src/latexmkrc), under which \mathfrak (and,
%% for safety, \mathcal/\overline/\widehat) are NOT brace-transparent, so an
%% unbraced subscript/superscript use like `_\fp` fails with
%% "! Missing { inserted." / "\__um_group_begin:". Every macro below whose
%% expansion starts with one of those commands is therefore given an extra
%% {...} layer in its own definition, so `_\fp`, `^\Qbar`, etc. keep working
%% unbraced at any call site. \DeclareMathOperator macros are unaffected
%% (\mathop already groups their content) and are ported unchanged.
\newcommand{\Q}{\mathbb{Q}}
\newcommand{\Z}{\mathbb{Z}}
\newcommand{\C}{\mathbb{C}}
\newcommand{\F}{\mathbb{F}}
\newcommand{\Qbar}{{\overline{\Q}}}
\newcommand{\Sha}{\mathrm{III}}
\newcommand{\fp}{{\mathfrak{p}}}
\newcommand{\fq}{{\mathfrak{q}}}
\newcommand{\ff}{{\mathfrak{f}}}
\newcommand{\OK}{{\mathcal{O}_K}}
\newcommand{\Lp}{L_p}
\newcommand{\Gm}{{\widehat{\mathbb{G}}_m}}
\DeclareMathOperator{\rank}{rank}
\DeclareMathOperator{\ord}{ord}
\DeclareMathOperator{\Reg}{Reg}
\DeclareMathOperator{\Sel}{Sel}
\DeclareMathOperator{\corank}{corank}
\DeclareMathOperator{\Gal}{Gal}
\DeclareMathOperator{\Tr}{Tr}
\DeclareMathOperator{\Nm}{N}
\DeclareMathOperator{\car}{char}
